% Originally: content/week-11.tex, \section{Solutions}

\section{Solutions}
\label{sec:flux_solutions}

\begin{exercisesolution}[Solution of \cref{ex:11.5}]
\textbf{(1)} Writing $V_\alpha^i(x) = \lvert x\rvert^\alpha x_i$,
\[
\partial_i\big(\lvert x\rvert^\alpha x_i\big) = \alpha\lvert x\rvert^{\alpha-2}x_i^2
  + \lvert x\rvert^\alpha \qquad \text{(no sum),}
\]
so, summing over $i=1,\ldots,n$,
\[
\divg V_\alpha(x) = \sum_{i=1}^n \Big(\alpha\lvert x\rvert^{\alpha-2}x_i^2 + \lvert
  x\rvert^\alpha\Big) = \alpha\lvert x\rvert^{\alpha-2}\lvert x\rvert^2 + n\lvert x\rvert^\alpha
  = (\alpha+n)\lvert x\rvert^\alpha.
\]
Since $\lvert x\rvert^\alpha \neq 0$ for $x\neq0$, this vanishes identically exactly when
$\alpha = -n$.

\textbf{(2)} For $\alpha=0$, $V_0(x)=x$ and, by part (1), $\divg V_0 \equiv n$. Applying
\cref{thm:gauss_divergence} on $B_r$,
\[\int_{B_r} \divg V_0\,dx = n\,\vol_n(B_r).\]
On $\partial B_r$ the outward unit normal is $\nu(x) = x/r$, so $V_0(x)\cdot\nu(x) = x\cdot
x/r = \lvert x\rvert^2/r = r$ (using $\lvert x\rvert=r$ on $\partial B_r$), giving
\[\int_{\partial B_r} V_0\cdot\nu\,d\vol_{n-1} = r\,\vol_{n-1}(\partial B_r).\]
Equating the two expressions of $\int_{B_r}\divg V_0\,dx = \int_{\partial B_r}V_0\cdot\nu\,
d\vol_{n-1}$ gives $r\,\vol_{n-1}(\partial B_r) = n\,\vol_n(B_r)$.

\textbf{(3)} By part (1), $\divg V_{-n}(x) = 0$ for every $x\neq0$, so $\int_{B_1}\divg
V_{-n}(x)\,dx = 0$ trivially (the integrand vanishes on $B_1\setminus\{0\}$, a set of full
measure in $B_1$). On $\partial B_1$, $\nu(x)=x$ and $\lvert x\rvert=1$, so $V_{-n}(x) =
\lvert x\rvert^{-n}x = x$, giving
\[\int_{\partial B_1} V_{-n}\cdot\nu\,d\vol_{n-1} = \int_{\partial B_1} x\cdot x\,d\vol_{n-1}
  = \int_{\partial B_1} 1\,d\vol_{n-1} = \vol_{n-1}(\partial B_1) > 0.\]
This does not contradict \cref{thm:gauss_divergence}: the theorem requires $F$ to be $C^1$ on
all of $\overline{B_1}$, but $V_{-n}(x) = \lvert x\rvert^{-n}x$ is not even defined, let alone
continuously differentiable, at $x=0 \in \overline{B_1}$. Its hypotheses simply do not hold on
this domain.
\end{exercisesolution}

% Verified: solved independently, all three parts match Sol11 exactly.

\begin{exercisesolution}[Solution to \cref{ex:ai_flux_cubic_sphere}]
% Generator: Claude Opus 5 (High)
\textbf{1.} On the unit sphere the exterior unit normal is $\nu(x,y,z) = (x,y,z)$, so the
integrand of \cref{def:flux} is $\langle F,\nu\rangle = x^4+y^4+z^4$. In the spherical
parametrization that becomes
$\sin^4\theta\cos^4\varphi + \sin^4\theta\sin^4\varphi + \cos^4\theta$, multiplied by the surface
element $\sin\theta$: a fourth-power trigonometric integral in two variables, where the
divergence theorem offers a radial integral in one.

\textbf{2.} Writing $r$ for the distance to the origin,
$\divg F = 3x^2+3y^2+3z^2 = 3r^2$, which is radial.

\textbf{3.} Because the integrand is radial, \cref{rem:polar_as_shells} applies with $n = 3$ and
$C_3 = 4\pi$:
\[\int_{\partial\Omega}\langle F,\nu\rangle\,d\vol_2 = \int_{B_1(0)} 3r^2\,d\vol_3
  = 4\pi\int_0^1 3r^2\cdot r^2\,dr = 12\pi\int_0^1 r^4\,dr = \frac{12\pi}{5}.\]

\begin{remark}[The direct route, for once, is also available]
\label{rem:cubic_flux_direct_check}
The surface integral abandoned in part 1 is not actually intractable, and working it out is a
worthwhile check. By the symmetry of the sphere under permuting the coordinates, the three terms
of $x^4+y^4+z^4$ contribute equally, so
\[\int_{\mathbb{S}^2}\big(x^4+y^4+z^4\big)\,d\vol_2 = 3\int_{\mathbb{S}^2} z^4\,d\vol_2
  = 3\int_0^{2\pi}\!\!\int_0^\pi \cos^4\theta\,\sin\theta\,d\theta\,d\varphi
  = 6\pi\left[-\frac{\cos^5\theta}{5}\right]_0^\pi = \frac{12\pi}{5},\]
in agreement. Note that what made this bearable was a symmetry argument, applied \emph{before}
any trigonometry: expanding all three fourth powers first is the version that is unpleasant.
\end{remark}
\end{exercisesolution}

\begin{exercisesolution}[Solution to \cref{ex:ai_capping_a_paraboloid}]
% Generator: Claude Opus 5 (High)
\textbf{1.} Take $\Omega := \{(x,y,z) : x^2+y^2 < 1,\ 0 < z < 1-x^2-y^2\}$, the solid region
between the paraboloid and the plane $z = 0$. Its boundary is $S$ together with the flat disk
\[D := \{(x,y,0) : x^2+y^2 \leq 1\},\]
whose exterior unit normal is $\nu \equiv -e_3$, since $\Omega$ lies above it. On $S$ the
exterior normal is the upward-pointing one, which is the orientation the exercise already fixes,
so no sign has to be corrected later.

\textbf{2.} Here $\divg F = 1+1+1 = 3$, so the flux out of $\partial\Omega$ is three times a
volume. In cylindrical coordinates,
\[\vol_3(\Omega) = \int_0^{2\pi}\!\!\int_0^1\!\!\int_0^{1-r^2} r\,dz\,dr\,d\varphi
  = 2\pi\int_0^1 \big(r - r^3\big)\,dr = 2\pi\left(\frac12 - \frac14\right) = \frac{\pi}{2},\]
and therefore
\[\int_{\partial\Omega}\langle F,\nu\rangle\,d\vol_2 = 3\,\vol_3(\Omega) = \frac{3\pi}{2}.\]

\textbf{3.} Along $D$ the third coordinate is $z = 0$, so
$\langle F,\nu\rangle = \langle (x,y,0),\,-e_3\rangle = -z = 0$ at every point: the field is
tangent to the cap, and contributes nothing. Splitting the boundary integral,
\[\int_S \langle F,\nu\rangle\,d\vol_2
  = \int_{\partial\Omega}\langle F,\nu\rangle\,d\vol_2 - \int_D \langle F,\nu\rangle\,d\vol_2
  = \frac{3\pi}{2} - 0 = \frac{3\pi}{2}.\]

\begin{remark}[The cap vanished because of this field, not because it is a cap]
\label{rem:cap_contribution_is_not_generally_zero}
It is tempting to read part 3 as licence to ignore the lid. Change $F$ to
$\widetilde F(x,y,z) := (x,\,y,\,z+1)$ and nothing about the first two parts moves, since
$\divg\widetilde F = 3$ as before and the closed flux is still $3\pi/2$. But now
$\langle \widetilde F,\nu\rangle = -1$ on $D$, so the cap carries a flux of $-\pi$, and the flux
through the paraboloid is $3\pi/2 + \pi = 5\pi/2$ instead. The subtraction in part 3 is doing
real work; it merely had nothing to subtract.

\Cref{ex:11.1} is the same technique on a hemisphere, and the two are worth reading as a pair:
there the cap contributes $\pi/2$ and must be subtracted, here it contributes nothing. The
method is identical and only the arithmetic differs, which is the point.
\end{remark}
\end{exercisesolution}
